\section{Mathematische Vorüberlegungen}

Gegeben ist ein gleichschenkliges Dreieck mit frei wählbarer
Schenkellänge \(s>0\). Der Winkel zwischen den beiden gleich langen
Schenkeln sei mit \(\alpha\) bezeichnet. Die Basisseite des Dreiecks habe
die Länge \(2r\). An diese Basisseite ist ein Halbkreis mit Radius \(r\)
bündig angelagert.

Gesucht ist derjenige Winkel \(\alpha\), für den die Gesamtfläche aus
Dreiecksfläche und Halbkreisfläche maximal wird. Außerdem soll der
maximale Flächeninhalt in Abhängigkeit von \(s\) angegeben werden.

Im Folgenden wird \(\alpha\) zunächst im Bogenmaß betrachtet, also

\[
\alpha \in [0,\pi].
\]

Die Umrechnung in Grad erfolgt am Ende über

\[
\alpha_{\degree}
=
\alpha \cdot \frac{180^\circ}{\pi}.
\]

\subsection{Geometrische Beziehungen}

Da das Dreieck gleichschenklig ist, halbiert die Höhe auf die Basisseite
sowohl die Basisseite als auch den Winkel \(\alpha\). Dadurch entstehen
zwei kongruente rechtwinklige Dreiecke. In einem dieser rechtwinkligen
Teildreiecke gilt:

\[
\text{Hypotenuse} = s,
\]

\[
\text{halbe Basis} = r,
\]

\[
\text{Winkel am oberen Eckpunkt} = \frac{\alpha}{2}.
\]

Daraus folgt mit dem Sinus:

\[
\sin\left(\frac{\alpha}{2}\right)
=
\frac{r}{s}.
\]

Somit ergibt sich

\[
r
=
s\sin\left(\frac{\alpha}{2}\right).
\]

Die Höhe \(h\) des gleichschenkligen Dreiecks ergibt sich entsprechend mit
dem Kosinus:

\[
\cos\left(\frac{\alpha}{2}\right)
=
\frac{h}{s},
\]

also

\[
h
=
s\cos\left(\frac{\alpha}{2}\right).
\]

Damit sind sowohl die Dreiecksfläche als auch die Halbkreisfläche als
Funktion von \(s\) und \(\alpha\) beschreibbar.

\subsection{Flächeninhalt des Dreiecks}

Die Grundseite des Dreiecks hat die Länge \(2r\), die zugehörige Höhe ist
\(h\). Daher gilt für die Dreiecksfläche:

\[
A_\triangle
=
\frac{1}{2}\cdot 2r \cdot h.
\]

Das vereinfacht sich zu

\[
A_\triangle
=
r h.
\]

Einsetzen von

\[
r=s\sin\left(\frac{\alpha}{2}\right)
\]

und

\[
h=s\cos\left(\frac{\alpha}{2}\right)
\]

liefert

\[
A_\triangle(\alpha)
=
s\sin\left(\frac{\alpha}{2}\right)
\cdot
s\cos\left(\frac{\alpha}{2}\right).
\]

Also

\[
A_\triangle(\alpha)
=
s^2
\sin\left(\frac{\alpha}{2}\right)
\cos\left(\frac{\alpha}{2}\right).
\]

Mit der trigonometrischen Identität

\[
\sin(\alpha)
=
2\sin\left(\frac{\alpha}{2}\right)
\cos\left(\frac{\alpha}{2}\right)
\]

folgt

\[
\sin\left(\frac{\alpha}{2}\right)
\cos\left(\frac{\alpha}{2}\right)
=
\frac{1}{2}\sin(\alpha).
\]

Damit erhält man für die Dreiecksfläche

\[
A_\triangle(\alpha)
=
\frac{1}{2}s^2\sin(\alpha).
\]

Alternativ erhält man dieselbe Formel direkt aus der Standardformel für
die Dreiecksfläche bei zwei gegebenen Seiten und eingeschlossenem Winkel:

\[
A_\triangle
=
\frac{1}{2}ab\sin(\gamma).
\]

Da hier beide Schenkel die Länge \(s\) besitzen und der eingeschlossene
Winkel \(\alpha\) ist, folgt ebenfalls

\[
A_\triangle(\alpha)
=
\frac{1}{2}s^2\sin(\alpha).
\]

\subsection{Flächeninhalt des Halbkreises}

Der an die Basis angelagerte Halbkreis hat den Radius \(r\). Die Fläche
eines ganzen Kreises mit Radius \(r\) ist

\[
A_\circ
=
\pi r^2.
\]

Der Halbkreis besitzt daher den Flächeninhalt

\[
A_\text{halb}
=
\frac{1}{2}\pi r^2.
\]

Mit

\[
r
=
s\sin\left(\frac{\alpha}{2}\right)
\]

folgt

\[
A_\text{halb}(\alpha)
=
\frac{1}{2}\pi
\left(
s\sin\left(\frac{\alpha}{2}\right)
\right)^2.
\]

Somit ist

\[
A_\text{halb}(\alpha)
=
\frac{1}{2}\pi s^2
\sin^2\left(\frac{\alpha}{2}\right).
\]

\subsection{Gesamtfläche}

Die grau hinterlegte Gesamtfläche setzt sich aus Dreiecksfläche und
Halbkreisfläche zusammen:

\[
A(\alpha)
=
A_\triangle(\alpha)
+
A_\text{halb}(\alpha).
\]

Einsetzen der zuvor hergeleiteten Terme liefert

\[
A(\alpha)
=
\frac{1}{2}s^2\sin(\alpha)
+
\frac{1}{2}\pi s^2
\sin^2\left(\frac{\alpha}{2}\right).
\]

Damit ist die Gesamtfläche als Funktion des Winkels \(\alpha\) gegeben
durch

\[
\boxed{
A(\alpha)
=
\frac{s^2}{2}
\left(
\sin(\alpha)
+
\pi
\sin^2\left(\frac{\alpha}{2}\right)
\right)
}
\]

für

\[
s>0
\qquad \text{und} \qquad
\alpha\in[0,\pi].
\]

Für die spätere Optimierung ist es günstig, den Halbwinkelterm
umzuformen. Mit

\[
\sin^2\left(\frac{\alpha}{2}\right)
=
\frac{1-\cos(\alpha)}{2}
\]

erhält man

\[
A(\alpha)
=
\frac{1}{2}s^2\sin(\alpha)
+
\frac{1}{2}\pi s^2
\cdot
\frac{1-\cos(\alpha)}{2}.
\]

Also

\[
A(\alpha)
=
\frac{1}{2}s^2\sin(\alpha)
+
\frac{1}{4}\pi s^2
\left(
1-\cos(\alpha)
\right).
\]

Man kann \(s^2\) ausklammern:

\[
A(\alpha)
=
s^2
\left(
\frac{1}{2}\sin(\alpha)
+
\frac{\pi}{4}
\left(
1-\cos(\alpha)
\right)
\right).
\]

Äquivalent dazu:

\[
\boxed{
A(\alpha)
=
\frac{s^2}{4}
\left(
2\sin(\alpha)
+
\pi
\left(
1-\cos(\alpha)
\right)
\right)
}
\]

Diese Darstellung ist besonders geeignet, um nach \(\alpha\) abzuleiten.

\subsection{Optimierung nach dem Winkel}

Da \(s>0\) fest vorgegeben ist, ist \(s^2\) eine positive Konstante.
Der Winkel, der \(A(\alpha)\) maximiert, hängt deshalb nicht von der
konkreten Größe von \(s\) ab.

Wir betrachten

\[
A(\alpha)
=
\frac{s^2}{4}
\left(
2\sin(\alpha)
+
\pi
\left(
1-\cos(\alpha)
\right)
\right).
\]

Die erste Ableitung nach \(\alpha\) ist

\[
A'(\alpha)
=
\frac{s^2}{4}
\left(
2\cos(\alpha)
+
\pi \sin(\alpha)
\right).
\]

Ein notwendiges Kriterium für ein inneres Extremum ist

\[
A'(\alpha)=0.
\]

Da \(s>0\) gilt, ist

\[
\frac{s^2}{4}>0.
\]

Daher ist die Gleichung

\[
A'(\alpha)=0
\]

äquivalent zu

\[
2\cos(\alpha)
+
\pi\sin(\alpha)
=
0.
\]

Umstellen liefert

\[
2\cos(\alpha)
=
-\pi\sin(\alpha).
\]

Für ein inneres Extremum im Bereich \(\alpha\in(0,\pi)\) ist
\(\sin(\alpha)>0\). Daher darf durch \(\sin(\alpha)\) dividiert werden:

\[
2\frac{\cos(\alpha)}{\sin(\alpha)}
=
-\pi.
\]

Mit

\[
\frac{\cos(\alpha)}{\sin(\alpha)}
=
\frac{1}{\tan(\alpha)}
\]

folgt

\[
\frac{2}{\tan(\alpha)}
=
-\pi.
\]

Damit ist

\[
\tan(\alpha)
=
-\frac{2}{\pi}.
\]

Da

\[
\alpha\in(0,\pi)
\]

gilt und \(\tan(\alpha)<0\) sein muss, liegt der Winkel im zweiten
Quadranten. Deshalb ist

\[
\boxed{
\alpha_\text{max}
=
\pi
-
\arctan\left(\frac{2}{\pi}\right)
}
\]

der relevante Kandidat für das Maximum.

Numerisch ergibt sich

\[
\alpha_\text{max}
\approx
2.574681067
\]

im Bogenmaß. In Grad ist das

\[
\alpha_\text{max}
\approx
2.574681067
\cdot
\frac{180^\circ}{\pi}.
\]

Also

\[
\boxed{
\alpha_\text{max}
\approx
147.52^\circ
}
\]

\subsection{Nachweis, dass es sich um ein Maximum handelt}

Zur Klassifikation des Extremums betrachten wir die zweite Ableitung.
Aus

\[
A'(\alpha)
=
\frac{s^2}{4}
\left(
2\cos(\alpha)
+
\pi\sin(\alpha)
\right)
\]

folgt

\[
A''(\alpha)
=
\frac{s^2}{4}
\left(
-2\sin(\alpha)
+
\pi\cos(\alpha)
\right).
\]

Für den kritischen Winkel

\[
\alpha_\text{max}
=
\pi
-
\arctan\left(\frac{2}{\pi}\right)
\]

gilt

\[
\sin(\alpha_\text{max})>0
\]

und

\[
\cos(\alpha_\text{max})<0.
\]

Damit sind beide Terme

\[
-2\sin(\alpha_\text{max})
\]

und

\[
\pi\cos(\alpha_\text{max})
\]

negativ. Also gilt

\[
A''(\alpha_\text{max})<0.
\]

Der kritische Winkel liefert daher ein lokales Maximum.

Zusätzlich betrachten wir die Randwerte. Für \(\alpha=0\) ist

\[
A(0)
=
\frac{s^2}{4}
\left(
2\sin(0)
+
\pi(1-\cos(0))
\right).
\]

Wegen

\[
\sin(0)=0
\qquad \text{und} \qquad
\cos(0)=1
\]

folgt

\[
A(0)=0.
\]

Für \(\alpha=\pi\) ist

\[
A(\pi)
=
\frac{s^2}{4}
\left(
2\sin(\pi)
+
\pi(1-\cos(\pi))
\right).
\]

Wegen

\[
\sin(\pi)=0
\qquad \text{und} \qquad
\cos(\pi)=-1
\]

folgt

\[
A(\pi)
=
\frac{s^2}{4}
\left(
\pi(1-(-1))
\right).
\]

Also

\[
A(\pi)
=
\frac{s^2}{4}\cdot 2\pi
=
\frac{\pi}{2}s^2.
\]

Da der innere Extremwert, wie im folgenden Abschnitt gezeigt wird,
größer als \(\frac{\pi}{2}s^2\) ist, handelt es sich um das globale
Maximum auf dem gesamten Intervall \([0,\pi]\).

\subsection{Berechnung des maximalen Flächeninhalts}

Nun wird der maximale Flächeninhalt berechnet. Dazu setzen wir

\[
\alpha_\text{max}
=
\pi
-
\arctan\left(\frac{2}{\pi}\right)
\]

in die Flächenformel ein.

Zur Vereinfachung definieren wir

\[
\beta
=
\arctan\left(\frac{2}{\pi}\right).
\]

Dann gilt

\[
\alpha_\text{max}
=
\pi-\beta.
\]

Aus

\[
\tan(\beta)
=
\frac{2}{\pi}
\]

können Sinus und Kosinus von \(\beta\) bestimmt werden. Wir betrachten ein
rechtwinkliges Dreieck mit Gegenkathete \(2\), Ankathete \(\pi\) und
Hypotenuse

\[
\sqrt{\pi^2+4}.
\]

Daraus folgt

\[
\sin(\beta)
=
\frac{2}{\sqrt{\pi^2+4}}
\]

und

\[
\cos(\beta)
=
\frac{\pi}{\sqrt{\pi^2+4}}.
\]

Da

\[
\alpha_\text{max}=\pi-\beta
\]

gilt,

\[
\sin(\alpha_\text{max})
=
\sin(\pi-\beta)
=
\sin(\beta)
=
\frac{2}{\sqrt{\pi^2+4}},
\]

und

\[
\cos(\alpha_\text{max})
=
\cos(\pi-\beta)
=
-\cos(\beta)
=
-\frac{\pi}{\sqrt{\pi^2+4}}.
\]

Wir verwenden die umgeformte Flächenformel

\[
A(\alpha)
=
\frac{s^2}{4}
\left(
2\sin(\alpha)
+
\pi
\left(
1-\cos(\alpha)
\right)
\right).
\]

Einsetzen von \(\alpha=\alpha_\text{max}\) liefert

\[
A_\text{max}(s)
=
\frac{s^2}{4}
\left(
2\cdot
\frac{2}{\sqrt{\pi^2+4}}
+
\pi
\left(
1-
\left(
-\frac{\pi}{\sqrt{\pi^2+4}}
\right)
\right)
\right).
\]

Das vereinfacht sich zu

\[
A_\text{max}(s)
=
\frac{s^2}{4}
\left(
\frac{4}{\sqrt{\pi^2+4}}
+
\pi
\left(
1+
\frac{\pi}{\sqrt{\pi^2+4}}
\right)
\right).
\]

Ausmultiplizieren ergibt

\[
A_\text{max}(s)
=
\frac{s^2}{4}
\left(
\frac{4}{\sqrt{\pi^2+4}}
+
\pi
+
\frac{\pi^2}{\sqrt{\pi^2+4}}
\right).
\]

Die beiden Brüche werden zusammengefasst:

\[
\frac{4}{\sqrt{\pi^2+4}}
+
\frac{\pi^2}{\sqrt{\pi^2+4}}
=
\frac{\pi^2+4}{\sqrt{\pi^2+4}}.
\]

Da

\[
\frac{\pi^2+4}{\sqrt{\pi^2+4}}
=
\sqrt{\pi^2+4},
\]

folgt

\[
A_\text{max}(s)
=
\frac{s^2}{4}
\left(
\pi
+
\sqrt{\pi^2+4}
\right).
\]

Damit lautet der maximale Flächeninhalt

\[
\boxed{
A_\text{max}(s)
=
\frac{\pi+\sqrt{\pi^2+4}}{4}
s^2
}
\]

Numerisch ist

\[
\frac{\pi+\sqrt{\pi^2+4}}{4}
\approx
1.716446108.
\]

Daher kann man auch schreiben:

\[
\boxed{
A_\text{max}(s)
\approx
1.716446108\,s^2
}
\]

\subsection{Zusammenfassung der mathematischen Ergebnisse}

Der Winkel, bei dem die Gesamtfläche maximal wird, ist

\[
\boxed{
\alpha_\text{max}
=
\pi
-
\arctan\left(\frac{2}{\pi}\right)
}
\]

beziehungsweise

\[
\boxed{
\alpha_\text{max}
\approx
147.52^\circ.
}
\]

Der maximale Flächeninhalt in Abhängigkeit von der Schenkellänge \(s\) ist

\[
\boxed{
A_\text{max}(s)
=
\frac{\pi+\sqrt{\pi^2+4}}{4}
s^2.
}
\]

Auffällig ist, dass der optimale Winkel \(\alpha_\text{max}\) unabhängig
von \(s\) ist. Die Schenkellänge \(s\) beeinflusst also nicht die Lage des
Maximums, sondern nur die Größe des maximalen Flächeninhalts. Da der
Flächeninhalt proportional zu \(s^2\) ist, skaliert eine Verdopplung von
\(s\) den maximalen Flächeninhalt um den Faktor \(4\).